\subsection{Strict $AB$-union and frontier module}

\paragraph{Exact interface.}
The body fixes
\[
 P_A=V_4+a(V_3-V_4),
 \qquad P_B=V_4+b(V_5-V_4),
\]
and the exact geometric family
\[
 \mathcal R_{AB}(a,b)=
 \bigcup\left\{
 T\cap H:
 \begin{array}{l}
 T\text{ is a closed unit equilateral triangle},\\
 V_4,P_A,P_B\in T
 \end{array}
 \right\}.
\]

\paragraph{Variables and feasible set.}

Let $e_A,e_B$ be unit vectors with angle $120^\circ$, put
\[
 O=0,\qquad A=ae_A,\qquad B=be_B,
 \qquad
 C=\{ue_A+ve_B:u,v\ge0\},
\]
identifying $O,A,B$ with the images of $V_4,P_A,P_B$, respectively.  In
this chart the public feasible union is
\begin{equation}
 \mathcal R_{AB}(a,b)=
 H\cap C\cap
 \bigcup\left\{
 T:
 \begin{array}{l}
 T\text{ is a closed unit equilateral triangle},\\
 O,A,B\in T
 \end{array}
 \right\}.                                   \label{eq:ab-union}
\end{equation}
The finite description in \zcref{thm:cert-caliper} applies to the
four-point set $\{O,A,B,x\}$ and makes membership in \eqref{eq:ab-union} an
exact finite semialgebraic condition of degree at most four.

The strict supercritical branch admits a simpler closed form.  In the next
statement write $x=ue_A+ve_B$ and
$\lVert x\rVert^2=u^2+v^2-uv$.

\paragraph{Objective.}
Determine the feasible union $\mathcal R_{AB}(a,b)$ exactly and identify the
cyclic order of its two circular and two linear non-axis frontier pieces.

\paragraph{Exact result.}

\begin{theorem}[Strict supercritical $AB$-union]
\label{thm:strict-ab-union}
Assume
\[
 a+b>1,\qquad \rho=a^2+ab+b^2<1,
\]
put $h=\sqrt3/2$, $D=\sqrt{4\rho-3}$, and define
\begin{align*}
 \alpha&=\frac{h(a+2b-aD)}{2\rho},&
 \beta&=\frac{h(a-b+(a+b)D)}{2\rho},\\
 \gamma&=\frac{h(-a+b+(a+b)D)}{2\rho},&
 \delta&=\frac{h(2a+b-bD)}{2\rho}.
\end{align*}
Then all four coefficients are positive and
\begin{equation}
 \mathcal R_{AB}(a,b)=H\cap\bigl(
 (C\cap\mathcal D_A\cap\mathcal H_A)
 \cup(C\cap\mathcal D_B\cap\mathcal H_B)\bigr),
 \label{eq:strict-ab-identity}
\end{equation}
where
\begin{align*}
 \mathcal D_A&=\{(u,v):(u-a)^2+v^2-(u-a)v\le1\},\\
 \mathcal D_B&=\{(u,v):u^2+(v-b)^2-u(v-b)\le1\},\\
 \mathcal H_A&=\{(u,v):\alpha(u-a)+\beta v\le0\},\\
 \mathcal H_B&=\{(u,v):\gamma u+\delta(v-b)\le0\}.
\end{align*}
Its non-axis frontier consists, in counterclockwise order, of the
$A$-circle, the $A$-line, the $B$-line, and the $B$-circle.
\end{theorem}

\paragraph{Solution.}

\begin{figure}[!htbp]
\centering
\begin{minipage}[t]{.48\linewidth}
\centering
\resizebox{.96\linewidth}{!}{%
\begin{tikzpicture}[x=1.55cm,y=1.55cm,>=Latex]
  \node[panellabel,anchor=west] at (-1.05,2.18)
    {(a) exact strict-supercritical $AB$ union};
  \coordinate (O) at (0,0);
  \coordinate (A) at (2.25,0);
  \coordinate (B) at (-1.125,1.949);
  \coordinate (P1) at (1.78,.33);
  \coordinate (Qm) at (1.35,.66);
  \coordinate (Q0) at (.57,1.10);
  \coordinate (Qp) at (-.08,1.46);
  \coordinate (P4) at (-.62,1.78);
  \draw[black!45,line width=.75pt] (O)--(2.55,0);
  \draw[black!45,line width=.75pt] (O)--(-1.28,2.217);

  \fill[VertexOrange!13]
    (O)--(A).. controls (2.08,.10) and (1.93,.23)..(P1)--(Qm)--(Q0)--cycle;
  \fill[CenterBlue!12]
    (O)--(Q0)--(Qp)--(P4).. controls (-.82,1.88) and (-1.00,1.92)..(B)--cycle;

  \draw[DemandRed,line width=1pt]
    (A).. controls (2.08,.10) and (1.93,.23)..(P1);
  \draw[VertexOrange,line width=1pt] (P1)--(Q0);
  \draw[CenterBlue,line width=1pt] (Q0)--(P4);
  \draw[GapPurple,line width=1pt]
    (P4).. controls (-.82,1.88) and (-1.00,1.92)..(B);

  \fill (A) circle (.018) node[below=2pt,figlabel] {$A=ae_A$};
  \fill (B) circle (.018) node[left=2pt,figlabel] {$B=be_B$};
  \fill (O) circle (.018) node[below left=1pt,figlabel] {$O$};
  \node[figlabel,text=DemandRed,rotate=28] at (2.06,.25) {$A$-circle};
  \node[figlabel,text=VertexOrange,rotate=30] at (1.30,.72) {$A$-line};
  \node[figlabel,text=CenterBlue,rotate=31] at (.12,1.44) {$B$-line};
  \node[figlabel,text=GapPurple,rotate=21] at (-.78,1.93) {$B$-circle};
  \node[figlabel,align=center] at (.45,-.55)
    {$\mathcal R(a,b)=(C\cap\mathcal D_A\cap\mathcal H_A)
      \cup(C\cap\mathcal D_B\cap\mathcal H_B)$};
\end{tikzpicture}%
}
\end{minipage}\hfill
\begin{minipage}[t]{.48\linewidth}
\centering
\resizebox{.96\linewidth}{!}{%
\begin{tikzpicture}[x=1.55cm,y=1.55cm,>=Latex]
  \node[panellabel,anchor=west] at (-1.05,2.18)
    {(b) the three asymmetric frontier witnesses};
  \coordinate (O) at (0,0);
  \coordinate (A) at (2.25,0);
  \coordinate (B) at (-1.125,1.949);
  \coordinate (P1) at (1.78,.33);
  \coordinate (Qm) at (1.35,.66);
  \coordinate (Q0) at (.57,1.10);
  \coordinate (Qp) at (-.08,1.46);
  \coordinate (P4) at (-.62,1.78);
  \draw[black!45,line width=.75pt] (O)--(2.55,0);
  \draw[black!45,line width=.75pt] (O)--(-1.28,2.217);
  \draw[black!55,line width=.75pt]
    (A).. controls (2.08,.10) and (1.93,.23)..(P1)--(Q0)--(P4)
    .. controls (-.82,1.88) and (-1.00,1.92)..(B);
  \fill[DemandRed] (Qm) circle (.035) node[below right=2pt,figlabel] {$Q_-$};
  \fill[DemandRed] (Q0) circle (.035) node[right=2pt,figlabel] {$Q_0$};
  \fill[DemandRed] (Qp) circle (.035) node[above right=2pt,figlabel] {$Q_+$};

  \draw[black!35,dashed] (Qm) circle (.76);
  \draw[black!35,dashed] (Qp) circle (.76);
  \draw[flow] (Qm)--($(Qm)!1.42!(O)$);
  \draw[flow] (Q0)--($(Q0)!1.42!(O)$);
  \draw[flow] (Qp)--($(Qp)!1.42!(O)$);
  \node[figlabel,atlasbox,align=left,anchor=west] at (.80,1.57)
    {frontier excludes the actual supercritical role;\\
     moving unit-circle signs exclude the two adjacent roles;\\
     vertex distances exclude the three nonlocal roles};
  \node[figlabel,align=center] at (.45,-.55)
    {$Q_-,Q_0,Q_+\in\operatorname{int}(H)
      \setminus\bigcup_{i=0}^{5}U_i$};
\end{tikzpicture}%
}
\end{minipage}
\caption{The strict-supercritical $AB$ frontier and the three asymmetric
witnesses.  The source-conditioned union has, in counterclockwise order, an
$A$-circle arc, an $A$-line segment, a $B$-line segment, and a $B$-circle
arc.  The points $Q_-,Q_0,Q_+$ are selected on this exact frontier and are
then excluded from the six actual V roles by frontier, moving-circle, and
vertex-distance arguments.  The drawing is schematic; the labelled frontier
order and witness roles are exact.}
\label{fig:strict-ab-frontier-witnesses}
\end{figure}


\begin{proof}
The identities
\[
 4\rho-3(a+b)^2=(a-b)^2,
 \qquad
 (a+b)^2D^2-(a-b)^2=4\rho((a+b)^2-1)
\]
give $0<D<1$ and positivity of the four coefficients.

We apply the caliper criterion of \zcref{thm:cert-caliper}.  For a
point $x=(u,v)$ in the relative interior of $C$ but outside
$\triangle OAB$, the four hull vertices occur in the order $O,B,x,A$.
The two cone-axis hull edges cannot certify a unit triangle: their support
sums, divided by $h$, are respectively
\[
 \max\{a,u\}+\max\{b,v-u\},\qquad
 \max\{b,v\}+\max\{a,u-v\},
\]
and are at least $a+b>1$.  It remains to test the $Ax$ and $Bx$ calipers.

For the $Ax$ caliper put $p=u-a$ and
$r^2=p^2+v^2-pv$.  Exact evaluation of the three supports gives
\begin{equation}
 \frac{G_A}{h}=
 \frac{\max\{r^2,-ap+(a+b)v\}+N_A}{r},
 \quad
 N_A=\begin{cases}
 0,&p\le0,\\
 ap,&0\le p\le v,\\
 (a+b)p-bv,&p\ge v.
 \end{cases}                                      \label{eq:Ax-caliper}
\end{equation}
No support label is omitted in these three sectors.  If $p\le0$,
$G_A\le h$ is equivalent to
\[
 r\le1,\qquad -ap+(a+b)v\le r.
\]
The second inequality is the half-plane condition in
\eqref{eq:strict-ab-identity}, because
\[
 r^2-(-ap+(a+b)v)^2
 =\frac{(cp+dv)(c'p+d'v)}{4\rho},
\]
where
\[
 c=a+2b-aD,\quad d=a-b+(a+b)D,
 \quad c'=a+2b+aD,\quad d'=a-b-(a+b)D,
\]
and $c,c',d>0>d'$.  Thus the low sector is exactly
$\mathcal D_A\cap\mathcal H_A$.  In the middle sector $0<p\le v$ one has
$r\le v$, whereas the numerator in \eqref{eq:Ax-caliper} is at least
$(a+b)v>r$, so no fit is possible.

In the high sector $p\ge v$, the inequalities $G_A\le h$ imply
\begin{equation}
 r^2+(a+b)p-bv\le r,\qquad bp+av\le r.          \label{eq:Ax-high}
\end{equation}
Write $x-A=ry(\theta)$, $0\le\theta\le60^\circ$, with cone coordinates
\[
 y(\theta)=te_A+we_B,\quad
 t=\frac2{\sqrt3}\cos(\theta-30^\circ),\quad
 w=\frac2{\sqrt3}\sin\theta.
\]
Then \eqref{eq:Ax-high} gives
\[
 r\le f(\theta)=1-(a+b)t+bw,\qquad S(\theta)=bt+aw\le1.
\]
Let $n_B$ be the unit vector with dual coordinates $(\gamma,\delta)$.
The identities
\[
 b\gamma+(a+b)\delta=h,\quad
 (a+b)\gamma+a\delta=\frac h2(1+D),\quad
 b\delta-a\gamma=\frac h2(1-D)
\]
show that the unique angle $\theta_B$ of this normal satisfies
$S(\theta_B)=1$ and $f(\theta_B)=(1-D)/2$.  The function $S$ has at most
one critical point, a maximum, so $S(\theta)\le1$ implies
$\theta\le\theta_B$.  On the feasible part $f,f'\ge0$; hence
$f(\theta)\cos(\theta_B-\theta)$ is increasing.  Consequently
\[
 \langle n_B,x-B\rangle
 \le a\gamma-b\delta+\frac h2(1-D)=0.
\]
The diameter bound also gives $\lVert x-B\rVert\le1$.  Thus every high
sector fit belongs to $\mathcal D_B\cap\mathcal H_B$.  Reflection exchanges
$a,u$ with $b,v$ and supplies the corresponding result for the $Bx$
caliper.  This proves \eqref{eq:strict-ab-identity} in the cone interior;
the axes follow by taking limits, since both sides are closed.

For completeness, the line--circle junctions are obtained by solving the
two displayed line and circle equations.  The line normals satisfy
\[
 \alpha^2+\alpha\beta+\beta^2=h^2,\qquad
 \gamma^2+\gamma\delta+\delta^2=h^2,
\]
and the line determinant is
\[
 \alpha\delta-\gamma\beta
 =\frac{3(1-D)(D+3)}{16\rho}>0.
\]
The successive cone slopes are
\[
 +\infty,\quad \frac2{1-D},\quad
 \frac{a(1-D)+2b}{2a+b(1-D)},\quad
 \frac{1-D}{2},\quad0;
\]
the two middle comparisons reduce to
$a(4-(1-D)^2)>0$ and $b(4-(1-D)^2)>0$.  This proves the asserted four-piece
order.
\end{proof}

\paragraph{Body consequence.}
This supplies the exact feasible-union and maximal-line-piece evaluation used
to define $\ell_-,\ell_+$ and $Q_0$ in the body, and is cited there as
\zcref{thm:strict-ab-union}.

\subsection{Handoff and nine-point frontier-forcing module}

\paragraph{Exact interface.}
The body interface is the exact chain and parameter domain in
\zcref{eq:reader-extreme-chain,eq:reader-ab-parameters,eq:reader-ab-domain},
the forced disk in \zcref{eq:reader-forced-disk}, the three frontier points
in \zcref{eq:reader-forced-asymmetric}, and the witness set in
\zcref{eq:reader-witness-set}.

\paragraph{Variables and feasible set.}
Assume $N_{\mathrm{gap}}=0$, $N_+=1$, and $(d,t)=(0,0)$; rotate the unique
supercritical role to index~$4$ and use the strict handoffs $\xi_i$ selected
in the body.  The feasible parameter pair is the strict component
$0<a,b<1$, $a+b>1$, $a^2+ab+b^2<1$ determined below.

\paragraph{Objective.}
Determine the common radial envelope and the two line--circle roots, then
prove that the resulting six radial points and three asymmetric frontier
points are excluded from every V triangle.

\paragraph{Exact result.}
The exact outputs are supplied by
\zcref{lem:ab-extreme-jump,lem:symmetric-core-witness,lem:fixed-line-signs,lem:asymmetric-core-witness}:
they give the strict endpoint inequalities, the forced centered disk, the
unique frontier roots, and $K_{\rm wit}(a,b)\subset U_C$.

\paragraph{Solution.}

\subsubsection*{Handoff stage}

\begin{lemma}[Exact-one handoff chain]
\label{lem:ab-extreme-jump}
Assume the six V triangles cover $\partial H$ and exactly one actual V triangle is
supercritical.  Rotate its index to $4$.  The strict handoffs supplied by
\zcref{prop:strict-handoffs} satisfy
\[
 \xi_3<\xi_2<\xi_1<\xi_0<\xi_5<\xi_4.                       \label{eq:extreme-chain}
\]
With
\[
 a=1-\xi_3,\qquad b=\xi_4,\qquad p=1-b,\qquad q=1-a,
\]
one has
\begin{equation}
 \begin{gathered}
 0<a,b<1,
 \qquad a=a_4<A_4,
 \qquad b=b_4<B_4,\\
 a+b>1,
 \qquad a^2+ab+b^2<1,
 \end{gathered}
 \label{eq:ab-strict-domain}
\end{equation}
and every selected V triangle obeys
\[
 a_i\ge p,\qquad b_i\ge q.
\]
\end{lemma}

\begin{proof}
Every selected V triangle other than $4$ is strictly nonsupercritical, so
$\xi_i<\xi_{i-1}$ for $i\ne4$, which is precisely
\eqref{eq:extreme-chain}.  In particular $\xi_3$ is the global minimum and
$\xi_4$ the global maximum.  Strictness of the selected handoffs gives
$a=1-\xi_3=a_4<A_4$ and $b=\xi_4=b_4<B_4$.  The two anchors of V
triangle $4$ lie in its open triangle;
their mutual squared distance is $a^2+ab+b^2$, strictly below the diameter
squared of the closed unit triangle.  This proves \eqref{eq:ab-strict-domain}.
Finally every $\xi_j\in[\xi_3,\xi_4]$, so
$1-\xi_{i-1}\ge1-\xi_4=p$ and $\xi_i\ge \xi_3=q$.
\end{proof}

\subsubsection*{Radial stage}

Put $h=\sqrt3/2$ and
\[
 s=p+q,\qquad m=\min(p,q),\qquad M=\max(p,q),
 \qquad \chi=s^4-s^2+pq.
\]
The strict domain \eqref{eq:ab-strict-domain} gives
\begin{equation}
 pq>(1-s)(2-s),\qquad m>1-s,\qquad s>1-m.         \label{eq:pq-domain}
\end{equation}
Let $c_*=c_{\max}(p,q)$ be the exact radial envelope of
\zcref{prop:exact-admissible-set}.  On the present subcritical
domain, equation \eqref{eq:radial-envelope} becomes
\begin{equation}
 c_*=\begin{cases}
 C_L(m),&\chi\le0,\\[1mm]
 \displaystyle\frac{2M}{1+\sqrt{4s^2-3}},&\chi>0,
 \end{cases}                                    \label{eq:core-cstar}
\end{equation}
where $C_L(m)$ is the unique root in $[\sqrt3/2,1]$ of
\[
 z^4-z^2+mz-m^2=0.
\]
At $\chi=0$ the two expressions agree and equal $s$.

\begin{lemma}[Direct radial forcing]
\label{lem:symmetric-core-witness}
One has
\begin{equation}
 0<c_*<1,\qquad c_*>1-m.                         \label{eq:cstar-bounds}
\end{equation}
Assume that every $U_i$ is $\Vdzero$ and that
$U_C,U_0,\ldots,U_5$ cover $H$.  Then, for every $i$,
\[
 P_i^{\mathrm{rad}}=(1-c_*)V_i\in U_C.
\]
Consequently
\begin{equation}
 \mathbb B(0,r_*)=\{x:\lVert x\rVert\le r_*\}
 \subseteq\conv\{P_0^{\mathrm{rad}},\ldots,P_5^{\mathrm{rad}}\}
 \subset U_C,
 \qquad r_*=h(1-c_*).                           \label{eq:forced-disk}
\end{equation}
\end{lemma}

\begin{proof}
On the $C_L$ branch, $C_L(m)<1$ because the defining polynomial is positive
at $1$; if $s<\sqrt3/2$, then $c_*\ge\sqrt3/2>s>1-m$, while if
$s\ge\sqrt3/2$ its value at $s$ is $\chi\le0$, so $c_*\ge s>1-m$.
On the other branch put $E=\sqrt{4s^2-3}$.  The selector gives
$sE>M-m$, whence $c_*<s<1$.  Moreover, with
\[
 A_*=\frac{2M}{1-m}-1,
\]
one has
\[
 A_*^2-E^2=
 \frac{4(1-s)\{(1-m)(1-2m)+m(2-m)(1-s)\}}{(1-m)^2}>0.
\]
Thus $A_*>E$ and $c_*>1-m$.  This proves
\eqref{eq:cstar-bounds}.

Fix $i$.  If $P_i^{\mathrm{rad}}\in U_i$, openness supplies $\varepsilon>0$ with
$c_*+\varepsilon<1$ such that
$(1-c_*-\varepsilon)V_i\in U_i$.  The same triangle contains its two selected
handoff anchors.  Since their incoming and forward-reach lower bounds dominate $p$ and
$q$, convexity and passage to the closure give a closed unit triangle
realizing $(p,q,c_*+\varepsilon)$.  This contradicts the definition of
$c_*=c_{\max}(p,q)$.  Hence $P_i^{\mathrm{rad}}\notin U_i$.

If $P_i^{\mathrm{rad}}$ belonged to $U_{i-1}$ or $U_{i+1}$, a small plane
ball inside that
open triangle would meet $r_i$ in a positive-length
interval.  This contradicts the $\Vdzero$ definition.  For the remaining
triangles put $r=1-c_*>0$.  Directly,
\[
 \lVert rV_i-V_{i\pm2}\rVert^2=1+r+r^2>1,
 \qquad
 \lVert rV_i-V_{i+3}\rVert=1+r>1.
\]
Since $V_j\in U_j$ and two points of an open unit triangle are at distance
strictly below $1$, these inequalities exclude $P_i^{\mathrm{rad}}$ from the
other three V triangles.  Thus no V triangle contains
$P_i^{\mathrm{rad}}$.  The point lies in
$\interior(H)$ by \eqref{eq:cstar-bounds}, so the assumed cover forces
$P_i^{\mathrm{rad}}\in U_C$.

The six points $P_i^{\mathrm{rad}}$ form a regular hexagon of circumradius
$1-c_*$.  Their
convex hull contains its centered incircle of radius $h(1-c_*)$, and convexity
of $U_C$ gives \eqref{eq:forced-disk}.
\end{proof}

\subsubsection*{Asymmetric frontier stage}

We next define the three asymmetric witnesses.  At $V_4$ use local
coordinates
\[
 \Phi(u,v):=V_4+u(V_5-V_4)+v(V_3-V_4),\qquad
 \mathsf q(u,v)=u^2+v^2-uv.
\]
In this chart the two public adjacent-handoff circles have centers
\[
 \begin{aligned}
 X_2(b)&=\Phi(E_-),& E_-&=(1-b,2-b),\\
 X_5(1-a)&=\Phi(E_+),& E_+&=(2-a,1-a),
 \end{aligned}
\]
and hence
\[
 \mathcal C_-=\partial\mathbb B(\Phi(E_-),1),
 \qquad
 \mathcal C_+=\partial\mathbb B(\Phi(E_+),1).
\]
The forward-reach lower bound is $b$ and the backward-reach lower bound is $a$.  Accordingly,
set
\begin{align*}
 \alpha&=\frac{h(b+2a-bD)}{2\rho},&
 \beta&=\frac{h(b-a+(a+b)D)}{2\rho},\\
 \gamma&=\frac{h(-b+a+(a+b)D)}{2\rho},&
 \delta&=\frac{h(2b+a-aD)}{2\rho},
\end{align*}
where now $\rho=a^2+ab+b^2$ and $D=\sqrt{4\rho-3}$.  The two line pieces of
\zcref{thm:strict-ab-union} are
\[
 L_-(\lambda)=(b-\beta\lambda,\alpha\lambda),\qquad
 L_+(\mu)=(\delta\mu,a-\gamma\mu).
\]
They meet at $\lambda_* =\mu_*=8h\rho/(3(D+3))$ in
\begin{equation}
 J=\left(\frac{2b+a-aD}{D+3},
          \frac{b+2a-bD}{D+3}\right).           \label{eq:line-junction}
\end{equation}
Thus the two public line pieces have the coordinate evaluations
\[
 \ell_-=\{\Phi(L_-(\lambda)):\lambda_*\le\lambda\le h^{-1}\},
 \qquad
 \ell_+=\{\Phi(L_+(\mu)):\mu_*\le\mu\le h^{-1}\}.
\]
The actual adjacent handoffs have local centers
\[
 Z_5(t)=(1+t,t),\qquad 1-a\le t\le b,
 \quad\text{and}\quad
 Z_3(y)=(y,1+y),\qquad 1-b\le y\le a.
\]
The restrictions of the two unit-circle equations are
\begin{align}
 g_-(\lambda)&=\frac34\lambda^2
 -3(\alpha+b\beta)\lambda+3b^2-3b+2,\notag\\
 g_+(\mu)&=\frac34\mu^2
 -3(\delta+a\gamma)\mu+3a^2-3a+2.              \label{eq:witness-quadratics}
\end{align}
Let $\lambda_-$ and $\mu_+$ be their first roots; explicitly,
\begin{align*}
 \lambda_-&=2(\alpha+b\beta)
 -\frac23\sqrt{9(\alpha+b\beta)^2-9b^2+9b-6},\\
 \mu_+&=2(\delta+a\gamma)
 -\frac23\sqrt{9(\delta+a\gamma)^2-9a^2+9a-6}.
\end{align*}

\begin{lemma}[Moving adjacent-triangle signs]
\label{lem:fixed-line-signs}
Throughout \eqref{eq:ab-strict-domain},
\begin{equation}
 \lambda_*<\lambda_-<h^{-1},\qquad
 \mu_*<\mu_+<h^{-1}.                            \label{eq:first-root-order}
\end{equation}
Moreover $J_u,J_v<1/2$.  The point $L_+(\mu_+)$ satisfies
\[
 (L_+(\mu_+))_u+(L_+(\mu_+))_v<3-2a,
\]
and the reflected bound holds for $L_-(\lambda_-)$.  Writing
$F(P,t)=\mathsf q(P-Z_5(t))-1$, one has, for every $t\in[1-a,b]$,
\begin{equation}
 F(L_+(\mu_+),t)\ge0,\qquad F(J,t)>0,\qquad
 F(L_-(\lambda_-),t)>0.                           \label{eq:moving-right-signs}
\end{equation}
Under the reflection
$(u,v,a,b)\leftrightarrow(v,u,b,a)$, the analogous signs hold for every
$y\in[1-b,a]$: the first-root point $L_-(\lambda_-)$ has nonnegative sign
against $Z_3(y)$, and the other two witnesses have strictly positive sign.
\end{lemma}

\begin{proof}
Put
\[
 U=a+b-1,\qquad z=a-b,\qquad
 D^2=z^2+6U+3U^2.
\]
Then $0<U<1/6$ and $z^2<\Omega:=1-6U-3U^2$.  We give the
fixed-line calculation for the circle centered at
$E_+=(2-a,1-a)$; reflection gives the calculation for $E_-$.  Put
\[
 F(P,t)=\mathsf q(P-(1+t,t))-1,\qquad t_0=1-a,
\]
so this circle is $F(P,t_0)=0$.  At the junction, exact expansion gives
\[
 (D+3)^2F(J,t_0)=3D(z^2+Uz+1-3U)+P_J,
\]
where
\[
 P_J=z^4+5z^2+6Uz^2+3U^2z^2+9Uz+3U+15U^2>0;
\]
indeed $z^2+Uz+1-3U\ge1-3U-U^2/4>0$, and
$5z^2+9Uz+15U^2$ is positive definite.  Thus $F(J,t_0)>0$.
Also
\[
 J_u+J_v=\frac{(1+U)(3-D)}{3+D}<1.
\]
Indeed this inequality is $3U<D(2+U)$, which follows from
$D^2\ge3U(2+U)$ and $(2+U)^3>3U$.  Since
\[
 \frac{\partial F}{\partial t}(P,t)=1+2t-P_u-P_v,
\]
$F(J,t)$ is strictly increasing for $t\ge0$, and therefore is positive
throughout $[1-a,b]$.

On the opposite line $L_-$, the derivative at the junction, after
multiplication by a positive factor, has at $t=1-a$ the numerator
\[
 N_0=D(9+3z^2+6Uz-9U^2)+P_0,
\]
where
\[
\begin{aligned}
 P_0={}&18U^3+6U^2z^2+63U^2+18Uz^2+18Uz+54U\\
      &+2z^4+9z^2+9.
\end{aligned}
\]
The coefficient of $D$ is at least $9-6U-9U^2>0$, while
$P_0\ge9+54U-18U>0$.  Hence $N_0>0$.  At the other endpoint $t=b$, the
corresponding numerator is
\[
 N_1=D(9+3z^2+6U-3U^2)+P_1,
\]
where
\[
\begin{aligned}
 P_1={}&9U^4-3U^3z+45U^3+9U^2z^2-6U^2z+72U^2\\
      &-Uz^3+21Uz^2+9Uz+45U+2z^4+9z^2+9.
\end{aligned}
\]
Its $D$-coefficient is positive and, using $|z|<1$,
\[
 P_1\ge9+45U-9U-U-6U^2-3U^3>0.
\]
The junction derivative is affine in $t$, so it is positive throughout
$[1-a,b]$.  Since the restriction of $F$ to $L_-$ is a convex quadratic,
is positive at the junction, and is increasing there, it has no intersection
on the opposite line side for any actual adjacent handoff.

On the correct line $L_+$, the value at its outer endpoint $h^{-1}$ has,
after multiplication by a positive factor, numerator
\[
 A_*=P_2-4DU(1+U+z),
\]
where
\[
\begin{aligned}
 P_2={}&3U^4+6U^3z+6U^3+4U^2z^2+12U^2z+12U^2\\
      &+2Uz^3+6Uz^2+2Uz+6U+z^4+2z^2-3.
\end{aligned}
\]
For fixed $U$, replace the odd powers of $z$ by $x=|z|$.  The resulting
polynomial $P_2^+(U,x)$ satisfies
\[
 \partial_xP_2^+
 =2(3U^3+4U^2x+6U^2+3Ux^2+6Ux+U+2x^3+2x)>0,
\]
and therefore its supremum for $z^2<\Omega$ occurs at $x=\sqrt\Omega$.
There
\[
 P_2^+(U,\sqrt\Omega)=4U(U+\sqrt\Omega-3)<0.
\]
As $1+U+z=2a>0$, this gives $A_*<0$.  The positive junction value and
negative endpoint value force the first root strictly between them, proving
the $L_+$ half of \eqref{eq:first-root-order}; reflection proves the other
half.

For the coordinate bounds, the exact identities
\[
 D^2-(3U-z)^2=6U(1+z-U),\qquad
 D^2-(3U+z)^2=6U(1-z-U)
\]
give $J_u,J_v<1/2$.  It remains to control the coordinate sum along $L_+$.
For $E=L_+(h^{-1})$, the sign of $3-2a-E_u-E_v$ is the sign of
\[
 D(6U+2z+6)+B(U,z),
\]
where
\[
\begin{aligned}
 B(U,z)={}&-9U^3-9U^2z-9U^2-3Uz^2-18Uz+3U\\
          &-3z^3+3z^2-3z+3.
\end{aligned}
\]
Here
\[
 B_z=-3(3z^2+2Uz+6U-2z+3U^2+1)<0,
\]
because the quadratic in parentheses has discriminant
$-32U^2-80U-8<0$.  Hence its minimum on $z^2<\Omega$ is at
$z=\sqrt\Omega$, where
$B(U,\sqrt\Omega)=6(1-3U-\sqrt\Omega)>0$ because
$(1-3U)^2-\Omega=12U^2$.  Thus $E_u+E_v<3-2a$.
Also $J_u+J_v<1<3-2a$, and the coordinate sum is affine on the
line segment, so the same strict bound holds at the first root.  Now
\[
 \frac{\partial F}{\partial t}(P,t)=1+2t-P_u-P_v.
\]
Since $F(L_+(\mu_+),1-a)=0$, the coordinate-sum bound makes its $t$-derivative
strictly positive for $t\ge1-a$.  Together with the no-opposite-line result,
this proves \eqref{eq:moving-right-signs}.  Reflection gives both the bound on
$L_-$ and the complete $Z_3(y)$ statement.
\end{proof}

Define the local and global witnesses by
\begin{equation}
 Q_-^{\rm loc}=L_-(\lambda_-),\qquad
 Q_0^{\rm loc}=J,\qquad
 Q_+^{\rm loc}=L_+(\mu_+),                       \label{eq:asym-witnesses}
\end{equation}
and
\[
 Q_-:=\Phi(Q_-^{\rm loc}),
 \qquad Q_0:=\Phi(Q_0^{\rm loc}),
 \qquad Q_+:=\Phi(Q_+^{\rm loc}).
\]
Thus $J$ is the unique meeting point of the two maximal line pieces, and the
first-root inequalities in \zcref{lem:fixed-line-signs} give exactly
\[
 \{Q_0\}=\ell_-\cap\ell_+,
 \qquad
 \{Q_-\}=\ell_-\cap\mathcal C_-,
 \qquad
 \{Q_+\}=\ell_+\cap\mathcal C_+.
\]

\begin{lemma}[Direct asymmetric forcing]
\label{lem:asymmetric-core-witness}
Assume that $U_C,U_0,\ldots,U_5$ cover $H$.  Then
\[
 Q_-,Q_0,Q_+
 \in\interior(H)\setminus\bigcup_{i=0}^5U_i,
\]
and consequently all three witnesses belong to $U_C$.
\end{lemma}

\begin{proof}
The root order and positivity of the line coordinates give
$0<u,v<1$ for all three points.  Each belongs to the local representation
of $\mathcal R_{AB}(a,b)$ obtained from \zcref{thm:strict-ab-union} by
interchanging the two cone coordinates, and hence
lies in a closed unit triangle containing $V_4$, so
$\mathsf q(u,v)\le1$; hence
$|u-v|<1$ and all three points lie in $\interior(H)$.

The closure of $U_4$ is one of the triangles represented by
$\mathcal R_{AB}(a,b)$, since it contains $V_4,X_3,X_4$.  All three witnesses lie
on its exact non-axis frontier.  If one belonged to the open triangle $U_4$,
a sufficiently small plane ball about it would lie both in $U_4$ and in the
interior of the corner cone, making the witness an interior point of
$\mathcal R_{AB}(a,b)$.  This contradiction excludes $U_4$.

For the triangles $U_0,U_1,U_2$, their contained vertices have local coordinates
\[
 \begin{array}{c|ccc}
 i&0&1&2\\ \hline
 V_i&(2,1)&(2,2)&(1,2).
 \end{array}
\]
For $P=(u,v)$,
\[
 \lVert P-(2,1)\rVert^2-1=\mathsf q(u,v)+2-3u,
 \quad
 \lVert P-(1,2)\rVert^2-1=\mathsf q(u,v)+2-3v.
\]
The inequalities $J_u,J_v<1/2$ and the line order give distance greater than
one from $V_0$ for $Q_-,Q_0$ and from $V_2$ for $Q_0,Q_+$.  For the remaining
pair, the first-root circle and the sum bound in
\zcref{lem:fixed-line-signs} give
\[
 \lVert Q_+-V_0\rVert^2-1
 =a(3-a-(Q_+)_u-(Q_+)_v)>a^2,
\]
and, by reflection, $\lVert Q_--V_2\rVert^2-1>b^2$.  All three points are
farther than one from $V_1=(2,2)$, because writing $u=1-\xi$, $v=1-\eta$
gives
\[
 \lVert P-V_1\rVert^2
 =1+\xi+\eta+\xi^2+\eta^2-\xi\eta>1.
\]
Since two points in one open unit equilateral triangle have distance strictly
less than one, these inequalities exclude $U_0,U_1,U_2$.

The actual handoff $X_5=Z_5(\xi_5)$ belongs to $U_5$, and
$1-a<\xi_5<b$ by \eqref{eq:extreme-chain}.  The three signs in
\eqref{eq:moving-right-signs} therefore put every witness at distance at
least one from $X_5$, excluding $U_5$.  Similarly, if $y=1-\xi_2$, then
$1-b<y<a$, $X_2=Z_3(y)\in U_3$, and the reflected signs exclude $U_3$.
Thus no V triangle contains any witness.  The assumed cover now forces all
three interior witnesses into $U_C$.
\end{proof}

The nine directly forced points are
$P_0^{\mathrm{rad}},\ldots,P_5^{\mathrm{rad}},Q_-,Q_0,Q_+$.  Under a
putative cover, \zcref{lem:symmetric-core-witness,lem:asymmetric-core-witness}
therefore force the compact set
\begin{equation}
 K_{\rm wit}(a,b)=
 \mathbb B(0,r_*)\cup\{Q_-,Q_0,Q_+\}
 \subset U_C.                                    \label{eq:forced-witness-set}
\end{equation}

\paragraph{Body consequence.}
This is the full calculation supplying the body containments
\zcref{eq:reader-forced-disk,eq:reader-forced-asymmetric,eq:reader-witness-set}.

\subsection{Newton reduction and four-overlap module}

\paragraph{Exact interface.}
Retain the public $K_{\rm wit}(a,b)$ from
\zcref{eq:reader-witness-set} and the exact support cap
\[
 \operatorname{Cap}(X,r)
 =\{n\in\mathbb S^1:\langle X,n\rangle+r\ge h\}.
\]

\paragraph{Variables and feasible set.}
Work on the strict $AB$ domain in \eqref{eq:ab-strict-domain} and its
nontrivial sheet $c_*>2/3$.  The Newton parameters are taken on the open
segments from the line junction $Q_0$ to the first circle roots $Q_-$ and
$Q_+$.

\paragraph{Objective.}
Replace the two independent first-root radicals by rational inner points and
verify the ray order, radius bounds, and four consecutive support-cap
intersections needed for the cyclic enclosure argument.

\paragraph{Exact result.}
The Newton containment is
\zcref{lem:technical-newton-reduction,eq:reader-newton-reduction}, and the
complete four-overlap conclusion is
\zcref{prop:technical-four-overlaps}.

\paragraph{Solution.}

For the nontrivial regime $c_*>2/3$, we first remove the square radicals in
the first-root witnesses.  Normalize
$\bar\alpha=\alpha/h$, and similarly for the other three coefficients, and
put
\[
 \xi=\lambda/h,\quad \zeta=\mu/h,\quad
 \xi_*=\zeta_*=\frac{8\rho}{3(D+3)}.
\]
Let $g_-,g_+$ be the rescaled forms of
\eqref{eq:witness-quadratics}, and take one Newton step at the junction:
\[
 \widehat\xi_-=\xi_*-\frac{g_-(\xi_*)}{g_-'(\xi_*)},\qquad
 \widehat\zeta_+=\zeta_*-\frac{g_+(\zeta_*)}{g_+'(\zeta_*)}.
\]
Define, in local coordinates,
\begin{align*}
 P_-&=\left(b-\frac34\bar\beta\widehat\xi_-,
          \frac34\bar\alpha\widehat\xi_-\right),\\
 P_0&=J,\\
 P_+&=\left(\frac34\bar\delta\widehat\zeta_+,
          a-\frac34\bar\gamma\widehat\zeta_+\right),
\end{align*}
and convert them to global vectors.

\begin{lemma}[Newton inner reduction]
\label{lem:technical-newton-reduction}
The points above have coordinates rational in $a,b,D$ and satisfy
\[
 P_-\in(Q_0,Q_-),\qquad P_0=Q_0,\qquad P_+\in(Q_0,Q_+).
\]
Consequently
\[
 \widehat K=\conv\bigl(\mathbb B(0,r_*)\cup\{P_-,P_0,P_+\}\bigr),
\]
and one has
\begin{equation}
 \widehat K\subseteq\conv(K_{\rm wit}),
 \qquad
 \Lambda(K_{\rm wit})\ge\Lambda(\widehat K).
 \label{eq:reader-newton-reduction}
\end{equation}
\end{lemma}

\begin{proof}
Each quadratic $g_-,g_+$ is strictly convex, positive, and decreasing at
the junction.  Its tangent zero therefore lies strictly between the
junction and its first root.  This gives the two open-segment inclusions.
All displayed coefficients and Newton quotients are rational in $a,b,D$,
and convexity gives the asserted containment.
\end{proof}

Put
\[
 P_*:=P_++\mathsf R P_-.
\]
For a disk $\mathbb B(X,r)$ let
\[
 \operatorname{Cap}(X,r)
 :=\{n\in\mathbb S^1:\langle X,n\rangle+r\ge h\}.
\]


\subsubsection*{Ray order and overlap verification}

\begin{proposition}[Technical four-overlap verification]
\label{prop:technical-four-overlaps}
Assume \eqref{eq:ab-strict-domain} and $c_*>2/3$.  Then
\[
 \lVert P_-\rVert,\lVert P_+\rVert>r_*,
\]
the rays $P_-,P_0,P_+,P_*,\mathsf R P_-$ occur in that cyclic order, and the
four consecutive
cap pairs
\[
 \operatorname{Cap}(P_-,2r_*)
 \leftrightarrow\operatorname{Cap}(P_0,2r_*)
 \leftrightarrow\operatorname{Cap}(P_+,2r_*)
 \leftrightarrow\operatorname{Cap}(P_*,r_*)
 \leftrightarrow\operatorname{Cap}(\mathsf R P_-,2r_*)
\]
overlap across their intervening short sectors.
\end{proposition}

\begin{proof}

Put $\tau=h-2r_*=h(2c_*-1)$.  The rays of
\begin{equation}
 P_-,\ P_0,\ P_+,\ P_*,\ \mathsf R P_-          \label{eq:cap-ray-order}
\end{equation}
occur in this cyclic order from $P_-$ to $\mathsf R P_-$.  Indeed the two
line distances
give
\[
 \frac{\operatorname{cross}(P_-,P_0)}{\lVert P_0-P_-\rVert}
 =\alpha(1-b)+\beta>0,
 \quad
 \frac{\operatorname{cross}(P_0,P_+)}{\lVert P_+-P_0\rVert}
 =\gamma+\delta(1-a)>0.
\]
If $P_-=(-X,-Y)$ and $P_+=(-P,-Q)$ in the centered cone basis, then
$0<X,Y,P,Q<1$, $X,Q>1/2$, and
\[
 \frac{\operatorname{cross}(P_+,\mathsf R P_-)}h
 =PX+(Q-P)Y>0.
\]
The last sign is immediate unless $Q<P$ and $X<Y$, in which case it is
larger than $Q-P/2>0$.  The ray of
$P_*=P_++\mathsf R P_-$ lies between those of $P_+$ and
$\mathsf R P_-$.  The same coordinates give
\begin{equation}
 \lVert P_-\rVert,\lVert P_+\rVert>h/2>r_*.
 \label{eq:outer-radii}
\end{equation}

For the two equal-radius overlaps it is enough to prove
\begin{equation}
 \frac{\operatorname{cross}(P_-,P_0)}{\lVert P_0-P_-\rVert}\ge\tau,
 \qquad
 \frac{\operatorname{cross}(P_0,P_+)}{\lVert P_+-P_0\rVert}\ge\tau.
 \label{eq:adjacent-overlaps}
\end{equation}
We now give the exact analytic verification.  By reflection take $a\le b$ and
set $m=1-b$, $M=1-a$, $U=a+b-1$, $s=m+M$, and $w=M-m$.  The smaller of the
two normalized left sides in \eqref{eq:adjacent-overlaps} is
\begin{equation}
 d=\frac{\mathcal A+U(2m-1)+D(\mathcal A+U)}{2\rho},
 \qquad \mathcal A=1-m+m^2.                      \label{eq:adjacent-d}
\end{equation}
On the $C_L$ sheet, write $F_L(z)=z^4-z^2+mz-m^2$.  This defining
polynomial evaluated at
$r=1-m/2+m^2/2$ is
\[
 F_L(r)=\frac{m^2(m-1)}{16}
 (m^5-3m^4+11m^3-17m^2+28m-12)>0.
\]
Indeed the polynomial in parentheses is increasing on $(0,1/2)$ and has
value $-33/32$ at $1/2$; its derivative is
$28-34m+m^2(5m^2-12m+33)>0$.  Since
$F_L(h)=-(m-h/2)^2\le0$ and the selected
root is unique, $2c_*-1\le\mathcal A$.  Put
\[
 X_0=\mathcal A+U,
 \qquad
 Y_0=2\rho\mathcal A-\mathcal A-U(2m-1).
\]
Then
\[
 d-\mathcal A=\frac{D(\mathcal A+U)-Y_0}{2\rho},
\]
and exact expansion gives
\[
\begin{aligned}
 Y_0={}&2(m^2-m+1)U^2+(2m^3-2m+3)U\\
      &+(m^2-m+1)(2m^2-2m+1)>0.
\end{aligned}
\]
Furthermore
\[
 D^2X_0^2-Y_0^2=4m\rho\Phi(U,m),
\]
where
\[
\begin{aligned}
 \Phi(U,m)={}&(1-m)(m^2-m+2)U^2\\
 &+(2-m^2)(m^2-m+1)U\\
 &-m(1-m)^2(m^2-m+1).
\end{aligned}
\]
In these variables
\[
 \chi=U^4-4U^3+5U^2-(m+2)U+m(1-m).
\]
Since $U^2(U^2-4U+5)>0$, the selector $\chi\le0$ implies
$U>m(1-m)/(m+2)$; $\Phi$ is increasing in $U$ and at that lower endpoint is
\[
 \frac{m^2(1-m)(m^4-2m^3+6m^2-6m+4)}{(m+2)^2}>0.
\]
The last quartic equals $(m^2-m)^2+5m^2-6m+4>0$.
Thus $d>\mathcal A\ge2c_*-1$.

On the other radial sheet let $E=\sqrt{4s^2-3}$.  Then
$0\le w<sE$ and $c_*=(s+w)/(1+E)$.  At fixed $U$, differentiation of
\eqref{eq:adjacent-d} can be checked without an implicit estimate.  Namely,
\[
 d=\frac{1+D}{2}-UB(D,w),\qquad
 B(D,w)=\frac{(1+U)(3+D)+w(1-D)}{D^2+3}.
\]
Here $D^2=w^2+6U+3U^2$, so $0\le D'=w/D\le1$, and
\[
 B_w=\frac{1-D}{D^2+3}\ge0.
\]
Writing $N=(1+U)(3+D)+w(1-D)$, direct differentiation gives
\[
 B_D=\frac{(1+U-w)(D^2+3)-2DN}{(D^2+3)^2}>-\frac{34}{27}.
\]
Indeed the first term in the numerator is positive, while
$1+U-w=2a>0$, $N<(7/6)4+1=17/3$, $D<1$, and $D^2+3\ge3$.  If $B_D<0$, then
$B_DD'\ge B_D$; otherwise $B_DD'\ge0$.  Hence
$B'=B_w+B_DD'>-34/27$, and $U<1/6$ yields
\[
 d'<\frac12+\frac16\frac{34}{27}=\frac{115}{162}<1,
\]
whereas
$(2c_*-1)'=2/(1+E)\ge1$.  Hence their difference decreases with $w$.
At the boundary $w=sE$, where $\chi=0$ and $c_*=s$, the preceding sheet
applies: this is an admissible point because $h<s<1$ and
$D^2=4s^4-12s+9<1$.  Indeed
$s^4-3s+2=(s-1)(s^3+s^2+s-2)<0$ on this interval; the second factor is
increasing and already equals $(7h-5)/4>0$ at $s=h$.  The preceding sheet gives
$d\ge2s-1$.  Therefore \eqref{eq:adjacent-overlaps} holds on both
sheets, and reflection supplies the other line.

It remains to check the two mixed overlaps.  The exact rational envelopes,
Gram reductions, and global positive-basis identities are recorded once in
\zcref{sec:strategy4-exact-certificate}.
\zcref{thm:paper-exact-mixed-certificate} gives precisely the overlaps for
$\mathbb B(P_+,2r_*)$ with $\mathbb B(P_*,r_*)$ and for
$\mathbb B(P_*,r_*)$ with $\mathbb B(\mathsf R P_-,2r_*)$.

Thus the four consecutive pairs of caps centered on the rays
\eqref{eq:cap-ray-order} overlap.
\end{proof}

\paragraph{Body consequence.}
These exact Newton and overlap data supply the nontrivial branch of the
support-arc solution below, and hence the body conclusion
\zcref{thm:reader-witness-enclosure}.
